\subsubsection{Physical stops, route patterns, and compact timetable index}
\label{sec:reduced-od-phase-two-timetable}

The second implementation phase compiles the domain timetable into
deterministic host arrays without constructing a time-expanded graph.  Scenario
stop identifiers remain the external OD and measurement geography.  A
\emph{physical stop place} is an explicit aggregation of one or more scenario
stops, typically platforms belonging to the same passenger interchange.  When
no mapping is supplied, the implementation uses an identity mapping.  A
nontrivial mapping must cover every scenario stop exactly, contain no unknown
stop, and be labeled either \texttt{authoritative} or
\texttt{reviewed\_generated}.  Place identifiers and member stop identifiers
are canonical and sorted; the representative coordinates are the arithmetic
mean of member coordinates.  The mapping policy, membership, and coordinates
enter its fingerprint.

A route pattern is defined by the tuple
\[
(\text{line identifier},\ \text{direction identifier},\
  \text{complete ordered physical-stop sequence}).
\]
Trips are grouped only when this tuple is identical.  Consequently, opposite
directions, express and local services, different lines, and repeated-stop
patterns remain distinct.  Platform variants may become equivalent only
through the explicit physical-stop mapping.  Every trip belongs to exactly one
route pattern, and both pattern and trip ordering are deterministic.

At this stage, a service-period class means a declared timetable
\texttt{service\_id}; the software does not infer calendar semantics or create
periods from OD departure bins.  Trips without a service identifier belong to
one reserved default class.  Exact scheduled arrival and departure times remain
in the timetable arrays, so after-midnight values above 24 hours are preserved.
Later range-routing phases may select trips by clock time without changing
these service classes.

The compact index contains sorted string dictionaries for scenario stops,
physical stops, trips, and lines.  Its numerical data consists of a trip
stop-time CSR pointer; aligned trip, scenario-stop, physical-stop, sequence,
arrival, and departure arrays for every stop-time record; and one line,
direction, route-pattern, and service-period index per trip.  Integer seconds
use 64-bit storage for timetable values, while bounded dictionary indices use
32-bit storage.  All arrays are owned, C-contiguous, and read-only.  Input
validation rejects duplicate or unknown identifiers, incomplete trips,
duplicate sequences, and nonmonotone times before publishing an artifact.
Artifact provenance independently fingerprints the source timetable, physical
stops, route patterns, service periods, and reduced-OD configuration.

The public Phase-2 CPU benchmark uses 5,000 trips with 15 stops per trip, or
75,000 stop-time records.  The complete normalization and array build takes
approximately 0.148 seconds, retains 2,520,008 array bytes, and increases the
observed process peak by approximately 25.5 MiB.  These machine-dependent
figures are not test thresholds, but confirm linear timetable storage and the
absence of time-expanded assignment state.

